\documentclass[10pt,UTF8,a4paper,twoside]{ctexart}

\usepackage{listings}
\usepackage{ctex}
\usepackage{multicol}
\usepackage{geometry}
\usepackage{graphicx}
\usepackage{subfigure}
\usepackage{fancyhdr}
\usepackage{appendix}
\usepackage{xcolor}
\usepackage{amsmath}
\usepackage{amssymb}

\geometry{a4paper,left=2cm,right=2cm,top=2cm,bottom=2cm}

% 用来设置附录中代码的样式
\lstset{
	language			=	C++,
    basicstyle          =   \ttfamily,          % 基本代码风格
    keywordstyle        =   \bfseries\color{blue},          % 关键字风格
    keywordstyle		=	[2] \color{teal},
    commentstyle        =   \rmfamily\itshape\color{red},  % 注释的风格，斜体
    stringstyle         =   \ttfamily\color{magenta},  % 字符串风格
    flexiblecolumns,                % 别问为什么，加上这个
    %numbers             =   left,   % 行号的位置在左边
    showspaces          =   false,  % 是否显示空格，显示了有点乱，所以不现实了
    %numberstyle         =   \zihao{-5}\ttfamily,    % 行号的样式，小五号，tt等宽字体
    showstringspaces    =   false,
    captionpos          =   t,      % 这段代码的名字所呈现的位置，t指的是top上面
    frame               =   single,   % 显示边框
    tabsize             =   4,
    breaklines          =   true,
    basewidth			=	0.5em,
}

\begin{document}

\begin{figure}[h]
\centering
\includegraphics{scut.jpg}
\end{figure}


%% temporary titles
% command to provide stretchy vertical space in proportion
\newcommand\nbvspace[1][3]{\vspace*{\stretch{#1}}}
% allow some slack to avoid under/overfull boxes
\newcommand\nbstretchyspace{\spaceskip0.5em plus 0.25em minus 0.25em}
% To improve spacing on titlepages
\newcommand{\nbtitlestretch}{\spaceskip0.6em}
\thispagestyle{empty}
\begin{center}
\bfseries
\nbvspace[1]
\Huge
\centering
{\nbtitlestretch\huge
ICPC TEMPLATE}

\begin{figure}[h]
\centering
\includegraphics[width=5cm,height=3cm]{ICPC.jpg}
\end{figure}

\nbvspace[1]
\normalsize


\nbvspace[1]
\small BY\\
\Large albertxwz\\[0.5em]
\footnotesize School of Computer Science \& Engineering,\\
South China University of Technology

\nbvspace[2]

%\includegraphics[width=1.5in]{./graphics/pic37}
\nbvspace[3]
\normalsize

Published in\\
\large
Nov 2020
\nbvspace[1]
\end{center}

\clearpage

\newpage
\begin{Large}
This template is a supplementary version.
\end{Large}
\mbox{}
\thispagestyle{empty}
\newpage
%*-----------------------------------------------------------------*
%|                             正文                       	     |
%*-----------------------------------------------------------------*
\begin{multicols}{2}
\columnseprule=0.5pt

\renewcommand \contentsname{Contents}
\tableofcontents
\thispagestyle{empty}

\clearpage


\setcounter{page}{1}
\pagestyle{headings}

\section{Standard Solution Template}
	\subsection{support bits/stdc++.h}
		\lstinputlisting[language=C++]{stdtemp/sol1.cpp}
	\subsection{unsupport bits/stdc++.h}
		\lstinputlisting[language=C++]{stdtemp/sol2.cpp}
	\subsection{Python Template}
		\lstinputlisting[language=python]{stdtemp/sol.py}
	
\clearpage

\section{Graph}
	\subsection{Network Flow}
		\subsubsection{Dinic}
			\lstinputlisting[language=C++]{graph/Dinic.cpp}
		\subsubsection{Edmonds-Karp}
			\lstinputlisting[language=C++]{graph/Edmonds-Karp.cpp}
	\subsection{DSU on Tree}
\clearpage

\section{Mathematics}
	\subsection{Number Theory}
		\subsubsection{Linear Inverse Modulo}
			\lstinputlisting[language=C++]{math/NumberTheory/LinInvMod.cpp}
		\subsubsection{M$\ddot{o}$bius Inversion}
		\subsubsection{Dujiao Sieve}
		\subsubsection{Min\_25 Sieve}
		\subsubsection{Lucas' Theorem}
			\begin{lstlisting}[language=C++]
long long Lucas(long long n, long long m, long long p) {
	if (m == 0) return 1;
	return (C(n % p, m % p, p) * Lucas(n / p, m / p, p)) % p;
}
			\end{lstlisting}
	\subsection{Karatsuba Multiply}
		\lstinputlisting[language=C++]{math/karatsuba.cpp}
	\subsection{Fast Fourier Transform}
		\lstinputlisting[language=C++]{math/FFT.cpp}
\clearpage

\section{Data Structure}
	\subsection{Treap}
	\subsection{Splay Tree}
	\subsection{Two-dimensional Segment Tree}
\clearpage

\section{String}
	\subsection{KMP}
		\lstinputlisting[language=C++]{string/KMP.cpp}
\clearpage

%*------------------------------------------------*
%|						附录						|
%*------------------------------------------------*

\begin{appendices}
	\section{Theorem}
		\large
		\subsection{Lucas' Theorem}
			对于质数p,有
			$$
				\tbinom{n}{m} \ mod \ p = \tbinom{\lfloor n/p \rfloor}{\lfloor m/p \rfloor}
				\cdot \tbinom{n \ mod \ p}{m \ mod \ p} \ mod \ p
			$$
		\subsection{Betty's Theorem}
			如果两个无理数$a,b$满足：
			$$ \frac{1}{a}+\frac{1}{b}=1$$
			那么对于两个集合$A,B$:
			$$A=\lbrace \lbrack na \rbrack \rbrace,B=\lbrace \lbrack nb \rbrack \rbrace$$
			有下面两个结论：
			$$A \bigcap B = \emptyset,A \bigcup B = \mathbb{N}^+$$
	\section{C++ STL: set}
		\large
		set与undered\_set的区别在于有无在内部存储时有无顺序。
		\subsection{Basic Method}
			\input{appendix/stlset/basicmethod.txt}
		\subsection{Advanced Method}
			注：必须导入algorithm头文件
			\input{appendix/stlset/advance.txt}
		
\end{appendices}

\end{multicols}
\end{document}